\newpage
\paragraph{Correct net cash can conceal incorrect gross cash flows.}
The snapshot worksheet leaves a period question unanswered: which legal-book
receipts and payments cross the group's boundary? An internal loan can produce
one company's investing payment and another company's financing receipt.
Adding the cash changes gives the right net amount, but retaining both flows
would misstate the group's gross activity. Matching activity labels is therefore
neither necessary nor sufficient to identify an internal transfer.

A separate constructed period illustrates the extension. At the start of
2025, P has cash 100, an investment in S of 100 and equity of 200; S has cash
100 and capital of 100. Continue to assume stable full ownership and control,
compatible policies and measurement bases, and investment cost equal to S's
capital. P then lends 30 to S, and S receives 20 from an external customer,
recognising income of 20. At year end P has cash 70 and a receivable of 30;
S has cash 150, a payable of 30 and income of 20. The investment and capital
remain unchanged. This constructed period has no interest, impairment,
acquisition difference or other transactions.

The opening worksheet removes the investment and capital of 100, leaving
group cash and equity of 200. The closing worksheet also removes the reciprocal
receivable and payable of 30, leaving cash 220 and equity including profit of
220. The period operation first obtains each legal book's own checked report.
It preserves all three original cash flows and names the internal payment and
receipt as a supplied pair. Removing that pair leaves an external operating
receipt of 20, no external payment and a net cash increase of 20. Profit of
20 also reconciles opening and closing group equity. The result is the same
whether both internal flows were labelled operating or the lender recorded
investing and the borrower financing.

A pair identifies the owner, flow identity and exact expected record on each
side, together with the supplied explanation connecting the transfers.
The dates and currencies must match; the amounts must be opposite and nonzero;
the owners must differ; and a flow cannot be removed twice. A changed record,
partial amount, duplicate reference or unexplained cash or equity movement
refuses. Opening and closing books must share the same period, chart,
ownership perimeter and compatibility assumptions. Each legal report retains
its own equity explanations and missing notes.

If the caller supplies no pair, the operation keeps the original gross flows.
It does not infer that two equal amounts describe the same transaction.
Likewise, an omitted closing reciprocal-balance adjustment leaves an extra
receivable and payable in the worksheet, even though equity still reconciles.
These cases remain visibly incomplete. By contrast, omitting the closing
investment adjustment breaks the declared equity bridge and refuses. The
distinction matters: reconciliation can expose contradictions without finding
every missing internal transaction.

The checked extension covers two account charts and both activity
classifications, with complete legal reports and group results preserved for
replay under \path{data/accounting-master/} in
\path{group-period-report-20261025}. It is a conditional period summary.
Actual source evidence, the completeness of paired transactions, indirect
operating cash flow, full notes and the wider group methods remain unresolved.

